\chapter{Results}

\section{Functional Verification of Conventional AES Architectures}

The conventional AES--128, AES--192, and AES--256 architectures were implemented and
verified using simulation testbenches. Each design was tested for both encryption and
decryption operations using standard NIST reference vectors. The results confirm full
functional correctness for all three key sizes.

\subsection{Testbench Output}

\begin{figure}[!htbp]
    \centering
    \includegraphics[width=0.95\linewidth]{figs/aes_pass_output.jpeg}
    \caption{Simulation console output showing PASS results for AES--128, AES--192, and AES--256.}
    \label{fig:aes_results_text}
\end{figure}
Figure~\ref{fig:aes_results_text} shows the console output generated during simulation.
For every key size, the encrypted ciphertext matches the expected value, and the decrypted
result successfully reconstructs the original plaintext.

\begin{itemize}
    \item AES--128: Encryption and decryption both reported \textbf{PASS}.
    \item AES--192: Encryption and decryption both reported \textbf{PASS}.
    \item AES--256: Encryption and decryption both reported \textbf{PASS}.
\end{itemize}

These results validate the correctness of the implemented AES round functions, key
expansion, and inverse transformations.

\FloatBarrier

\subsection{Waveform Verification}

To further confirm functional correctness, waveform analysis was conducted. 
Figure~\ref{fig:aes_waveform} illustrates the transitions of key internal signals during
encryption and decryption. The evolution of the output registers shows proper propagation
of intermediate states until the correct final values are reached.  
Additionally, the \texttt{passed\_tests} register increments successfully and
\texttt{failed\_tests} remains zero, indicating error-free simulation.

\begin{figure}[!htbp]
    \centering
    \includegraphics[width=0.95\linewidth]{figs/aes_waveform.jpeg}
    \caption{Waveform showing transitions of AES encryption and decryption signals.}
    \label{fig:aes_waveform}
\end{figure}

\FloatBarrier

\section{RTL Schematic Verification}


\begin{figure}[!htbp]
    \centering
    \includegraphics[width=\linewidth]{figs/aes_rtl_schematic.jpeg}
    \caption{RTL schematic showing instantiation and connectivity of AES modules for all key sizes.}
    \label{fig:aes_rtl}
\end{figure}

The RTL schematic generated by Vivado confirms that the structural composition of the AES
modules aligns with the architectural intent. As shown in
Figure~\ref{fig:aes_rtl}, the testbench instantiates three AES encryption and decryption
units corresponding to the 128, 192, and 256 bit key sizes and compares their outputs with
expected results.


\FloatBarrier

\section{Resource Utilization}

Synthesis was performed using Vivado to determine the hardware cost of the implemented
AES architectures. Table~\ref{tab:aes_baseline_comparison} compares the LUT and
flip-flop counts for the pipeline and FSM architectures at AES-128 before optimization.
The pipeline architecture, which instantiates one copy of the round logic per pipeline
stage, requires significantly more area than the FSM architecture, which reuses the same
round hardware across all rounds.

\begin{table}[!htbp]
    \centering
    \caption{Vivado synthesis resource utilization --- AES-128 baseline comparison}
    \label{tab:aes_baseline_comparison}
    \renewcommand{\arraystretch}{1.15}
    \begin{tabular}{lccc}
        \hline
        \textbf{Architecture} & \textbf{Slice LUTs} & \textbf{Flip-Flops} & \textbf{Notes} \\
        \hline
        Pipeline (AES-128 baseline) & 38,766 & 2,763 & EXP3/LN3 log-table MixColumns \\
        FSM (AES-128 baseline)      &  6,794 &   412 & EXP3/LN3 log-table MixColumns \\
        Pipeline (Phase~1 est.)    & $\approx$25k--30k & --- & xtime arithmetic (TBD Vivado) \\
        FSM (Phase~1 est.)         & $\approx$4k--5k  & --- & xtime arithmetic (TBD Vivado) \\
        \hline
    \end{tabular}
\end{table}

\FloatBarrier

The baseline figures confirm that the pipeline architecture consumes approximately
5.7$\times$ more LUTs than the FSM architecture for AES-128. Phase~1 is projected to
reduce pipeline LUTs by 25--35 percent and FSM LUTs by 25--30 percent through
elimination of the 512-entry EXP3 and LN3 tables from the MixColumns datapath. Actual
Phase~1 and Phase~2 synthesis numbers will be populated after Vivado runs complete.

\begin{table}[!htbp]
    \centering
    \caption{Vivado synthesis resource utilization --- conventional AES FSM
    architecture (AES-128), S-Box and key expansion module breakdown.}
    \label{tab:aes_resources}
    \begin{tabular}{lccc}
        \hline
        \textbf{Module} & \textbf{Slice LUTs} & \textbf{F7 Muxes} & \textbf{IOBs} \\
        \hline
        Key Expansion Round & 300 & 64 & 292 \\
        S-Box 0 & 72 & 16 & 0 \\
        S-Box 1 & 72 & 16 & 0 \\
        S-Box 2 & 72 & 16 & 0 \\
        S-Box 3 & 72 & 16 & 0 \\
        \hline
    \end{tabular}
\end{table}

\FloatBarrier

\subsection{LUT Reduction Across Architectures}

\begin{figure}[!htbp]
    \centering
    \includegraphics[width=0.85\linewidth]{figs/LUT.jpeg}
    \caption{LUT utilization comparison across pipeline and FSM architectures at AES-128
    baseline. The FSM architecture achieves approximately 82 percent lower LUT usage than
    the pipeline architecture at baseline. Phase~1 xtime replacement is projected to
    reduce both architectures further.}
    \label{fig:lut_comparison}
\end{figure}

\FloatBarrier

\section{Phase~1 and Phase~2 Results (Pending Vivado Runs)}

The following figures and tables will be populated after Vivado synthesis runs complete
for the optimized RTL. Placeholder elements are included here to indicate the intended
structure of the results section.

\subsection{Simulation Waveform Verification}

\begin{figure}[!htbp]
\centering
\fbox{\parbox{0.85\linewidth}{\centering\vspace{1.5cm}
\textbf{Pipeline Architecture --- Simulation Waveform}\\[0.4cm]
\textit{Insert Vivado/ModelSim waveform screenshot showing pipeline AES-128 encryption
signals after Phase~1 xtime replacement. Confirm output\_valid and ciphertext signals
match NIST test vectors.}
\vspace{1.5cm}}}
\caption{Pipeline AES architecture simulation waveform --- Phase~1 xtime replacement
(result pending).}
\label{fig:pipeline_waveform_ph1}
\end{figure}

\FloatBarrier

\begin{figure}[!htbp]
\centering
\fbox{\parbox{0.85\linewidth}{\centering\vspace{1.5cm}
\textbf{FSM Architecture --- Simulation Waveform}\\[0.4cm]
\textit{Insert Vivado/ModelSim waveform screenshot showing FSM AES-128 encryption
signals after Phase~2 clock enable guards. Confirm ready signal behavior and
ciphertext output match NIST test vectors. Verify CE guard freezes register during
IDLE cycles.}
\vspace{1.5cm}}}
\caption{FSM AES architecture simulation waveform --- Phase~2 clock enable guards
(result pending).}
\label{fig:fsm_waveform_ph2}
\end{figure}

\FloatBarrier

\subsection{Post-Optimization Synthesis Results}

\begin{table}[!htbp]
\centering
\caption{Pipeline architecture Vivado synthesis resource utilization --- Phase~1
xtime replacement (values TBD after Vivado run).}
\label{tab:pipeline_ph1_resources}
\renewcommand{\arraystretch}{1.15}
\begin{tabular}{lcccc}
    \hline
    \textbf{Phase} & \textbf{Slice LUTs} & \textbf{Flip-Flops} &
    \textbf{Delta LUTs} & \textbf{\% Reduction} \\
    \hline
    Baseline (EXP3/LN3) & 38,766 & 2,763 & --- & --- \\
    Phase~1 (xtime)     & TBD    & TBD   & TBD & TBD \\
    \hline
\end{tabular}
\end{table}

\FloatBarrier

\begin{table}[!htbp]
\centering
\caption{FSM architecture Vivado synthesis resource utilization --- Phase~1 and
Phase~2 optimizations (values TBD after Vivado run).}
\label{tab:fsm_ph2_resources}
\renewcommand{\arraystretch}{1.15}
\begin{tabular}{lcccc}
    \hline
    \textbf{Phase} & \textbf{Slice LUTs} & \textbf{Flip-Flops} &
    \textbf{Delta LUTs} & \textbf{\% Reduction} \\
    \hline
    Baseline (EXP3/LN3)       &  6,794 & 412 & --- & --- \\
    Phase~1 (xtime)           & TBD    & TBD & TBD & TBD \\
    Phase~2 (CE guards)       & TBD    & TBD & TBD & TBD \\
    \hline
\end{tabular}
\end{table}

\FloatBarrier

\section{Summary}

The results of the conventional AES implementation can be summarized as follows:

\begin{itemize}
    \item All three AES variants (128, 192, and 256 bit) were implemented successfully.
    \item Encryption and decryption outputs matched the expected values for all official test vectors.
    \item Waveform inspection confirms correct timing behavior and internal signal transitions.
    \item RTL schematics verify the structural correctness of the datapath and key schedule modules.
    \item Pipeline architecture (AES-128 baseline): 38,766 LUTs and 2,763 flip-flops.
    \item FSM architecture (AES-128 baseline): 6,794 LUTs and 412 flip-flops, representing
          an 82 percent area reduction relative to the pipeline architecture.
\end{itemize}

These outcomes form a solid foundation for the architectural optimizations and power aware
enhancements introduced in subsequent chapters.
